\section{问题四：给定径向收缩的移动边界模型}
\subsection{由干物质守恒推导材料运动}
问题四同时给出了半径变化和一套新的经验物性，不能仅把问题三的固定半径替换为末时半径。Brasiello等研究了食品干燥的非等温移动边界模型~\cite{brasiello2021}，将收缩与传热传质联系起来。本文直接利用附件2构造已知$R(t)$，按假设（2）取比例径向收缩速度
\begin{equation}\label{eq_velocity}
v(r,t)=\frac{\dot R(t)}{R(t)}r.
\end{equation}
外半径数据本身不能唯一确定内部运动，式\eqref{eq_velocity}是用于闭合速度场的假设。

设$s(r,t)$为单位当前体积的干物质质量，$j$为相对固体骨架的向外水质量通量。干物质不流失时，干物质与水质量分别满足
\begin{equation}\label{eq_solid_water}
\begin{aligned}
s_t+\frac1r\partial_r(rsv)&=0,\\
\partial_t(sC)+\frac1r\partial_r[r(sCv+j)]&=0,\qquad j=-sD(C,T)C_r.
\end{aligned}
\end{equation}
将第二式展开，减去第一式乘以$C$，得到
\begin{equation}\label{eq_material_C}
C_t+vC_r=\frac1{sr}\partial_r(rsDC_r).
\end{equation}
初始$s$均匀时，将式\eqref{eq_velocity}代入干物质方程，可得
\begin{equation}\label{eq_s_density}
\dot s+2\frac{\dot R}{R}s=0,\qquad
s(t)=s_0\left[\frac{R_0}{R(t)}\right]^2.
\end{equation}
因此$s$保持空间均匀，可从式\eqref{eq_material_C}的空间散度中约去。式\eqref{eq_solid_water}至式\eqref{eq_s_density}说明，纯收缩不会改变同一材料单元的水质量与干物质质量之比；对干基含水率直接添加$-C\nabla\cdot v$项会重复计算几何变化。

\subsection{材料坐标变换与控制方程}
定义固定区间上的材料坐标和状态
\begin{equation}\label{eq_xi}
\xi=\frac{r}{R(t)},\quad
U(\xi,t)=C(R(t)\xi,t),\quad
\Theta(\xi,t)=T(R(t)\xi,t),\qquad 0\le\xi\le1.
\end{equation}
由链式法则，$U_t=C_t+\xi\dot R C_r=C_t+vC_r$，且$C_r=U_\xi/R$。式\eqref{eq_material_C}于是变为
\begin{equation}\label{eq_shrink}
\boxed{\begin{aligned}
U_t&=\frac1{R(t)^2\xi}\partial_\xi[\xi D(U,\Theta)U_\xi],\\
B(U)\Theta_t&=\frac1{R(t)^2\xi}\partial_\xi[\xi k(U)\Theta_\xi].
\end{aligned}}
\end{equation}
材料运动与坐标网格同步，使显式对流项抵消。式\eqref{eq_shrink}的热方程是前述有效传热近似在材料坐标下的表达，不由水质量守恒自动推出，也未包含收缩机械功或完整水分携焓。

本问全部物性改为附录4经验式：
\begin{equation}\label{eq_prop4}
\begin{aligned}
\rho(U)&=760+90U,&c_p(U)&=1850+2150\frac{U}{1+U},\\
k(U)&=0.12+0.20\frac{U}{1+U},&&\\[-3pt]
D(U,\Theta)&=4.2\times10^{-4}\exp\!\left(-\frac{0.30}{U}\right)
\exp\!\left[-\frac{3850}{\Theta+273.15}\right].&&
\end{aligned}
\end{equation}
将式\eqref{eq_prop4}与$R(0)=0.02$ m、$U(\xi,0)=2.55$、$\Theta(\xi,0)=28$共同使用。中心仍取$U_\xi(0,t)=\Theta_\xi(0,t)=0$；外表面为
\begin{equation}\label{eq_shrink_boundary}
-\frac{D_s}{R}U_\xi(1,t)=h_m(U_s-C_a),\qquad
-\frac{k_s}{R}\Theta_\xi(1,t)=h(\Theta_s-T_a).
\end{equation}
式\eqref{eq_xi}至式\eqref{eq_shrink_boundary}共同保留了内部扩散的$R^{-2}$与边界梯度的$R^{-1}$，不能只改控制方程而遗漏边界尺度。

\subsection{固定材料网格上的求解与物理输出}
将问题一的$G_8$节点除以$R_0$得到固定材料节点$\xi_i$，令内部面$\xi_{i+1/2}=(\xi_i+\xi_{i+1})/2$，径向权重$w_i=(\xi_{i+1/2}^2-\xi_{i-1/2}^2)/2$，端点按0和1截取。记
\[
\widehat F^C_{i+1/2}=-\xi_{i+1/2}D_{i+1/2}
\frac{U_{i+1}-U_i}{\xi_{i+1}-\xi_i},\qquad
\widehat F^T_{i+1/2}=-\xi_{i+1/2}k_{i+1/2}
\frac{\Theta_{i+1}-\Theta_i}{\xi_{i+1}-\xi_i}.
\]
表面通量分别取$Rh_m(U_N-C_a)$与$Rh(\Theta_N-T_a)$。半离散方程为
\begin{equation}\label{eq_shrink_fvm}
R^2w_i\dot U_i=\widehat F^C_{i-1/2}-\widehat F^C_{i+1/2},\quad
R^2w_iB_i\dot\Theta_i=\widehat F^T_{i-1/2}-\widehat F^T_{i+1/2}.
\end{equation}
式\eqref{eq_shrink_fvm}使用2113个材料节点，BDF相对容差为$10^{-8}$，温度与含水率绝对容差分别为$10^{-8}$、$10^{-10}$；前4 h最大步长10 s，其后60 s。环境切换处分段，过程状态连续。以全部材料节点的最大含水率定位事件，并在细网格和时间收紧解上核验同一报告时刻。

题目要求的是当前物理距离。对每个$r_j\le R(t)$，先求$\xi_j=r_j/R(t)$，再对材料网格上的$U$作保形插值；固定$r_j>R(t)$的位置已在药材外，单元格留空。表面始终使用$\xi=1$的实际节点并单列，不能用固定的2 cm列表示移动表面。

\subsection{烘干时长与径向含水率}
生产解临界时刻为51.0871229810 h。程序在生产、空间加密和时间收紧解的最晚临界时刻之后增加2 s，再向上选择四位小时报告值，并复核三组未舍入最大值。最终得到
\begin{equation}\label{eq_q4_answer}
\boxed{t_{\mathrm{rep},4}=51.0877\ \mathrm h=183915.72\ \mathrm s.}
\end{equation}
式\eqref{eq_q4_answer}对应最大含水率为0.1499989333 kg/kg，最大值位于中心，当前半径为1.2000 cm，位于附件2的72 h覆盖范围内，无需半径外推。表\ref{tab_q4C}为题目要求的结果，图\ref{fig_q4}同时展示含水率历程与当前物理半径剖面。
\begin{table}[H]\centering\small
\caption{收缩模型药材烘干过程的水分浓度（kg/kg）}\label{tab_q4C}
\renewcommand{\arraystretch}{1.12}
\begin{tabular}{cccccccc}\toprule
 & \multicolumn{6}{c}{到药材中心的距离/cm} & \\
时间/h & 0 & 0.5 & 1.0 & 1.5 & 2.0 & 药材表面 & $R(t)$/cm\\\midrule
6 & 1.7188 & 1.5370 & 1.0221 & --- & --- & 0.4205 & 1.3740\\
12 & 0.7374 & 0.6544 & 0.4082 & --- & --- & 0.1669 & 1.2480\\
18 & 0.4086 & 0.3686 & 0.2397 & --- & --- & 0.0892 & 1.2140\\
24 & 0.2851 & 0.2610 & 0.1790 & --- & --- & 0.0673 & 1.2040\\
30 & 0.2263 & 0.2094 & 0.1493 & --- & --- & 0.0595 & 1.2010\\
36 & 0.1928 & 0.1796 & 0.1316 & --- & --- & 0.0560 & 1.2000\\
42 & 0.1712 & 0.1604 & 0.1200 & --- & --- & 0.0541 & 1.2000\\
48 & 0.1561 & 0.1468 & 0.1116 & --- & --- & 0.0530 & 1.2000\\
51.0877 & 0.1500 & 0.1413 & 0.1081 & --- & --- & 0.0526 & 1.2000\\
\bottomrule\end{tabular}
\par\vspace{3pt}\begin{minipage}{.96\linewidth}\footnotesize 注：---表示该固定距离已超出当时半径；药材表面对应当前$R(t)$。终时中心未舍入值为0.1499989333 kg/kg。\end{minipage}
\end{table}

\figpair{q4_C_history.pdf}{q4_C_profiles.pdf}{收缩模型的含水率变化。左图为中心和表面历程；右图以当前物理半径绘制剖面，曲线右端随收缩移动，始终对应真实表面。}{fig_q4}

表\ref{tab_q4C}各列含水率单位均为kg/kg，空白表示该固定距离已超出当时半径；附加的$R(t)$列用于明确表面位置。终时中心四位显示为0.1500，与问题三相同，应结合未舍入最大值理解严格达标条件。

第四问比第三问的报告时长减少6.3854 h，约为第三问的11.11\%。这个差异同时包含附录3变为附录4以及几何收缩两项改变，不能据此认定收缩单独使时长减少11.11\%。若要分离二者作用，需要另作相同物性下固定域与收缩域的对照，本文不将尚未完成的对照写成结论。

\subsection{收缩过程的干基收支}
按式\eqref{eq_s_density}，干物质总量$M_d=s_0\pi R_0^2L$保持不变，水质量满足
\begin{equation}\label{eq_mass_shrink}
\frac{M_w}{M_d}=2\int_0^1U\xi\dd\xi,\qquad
\frac{\dd}{\dd t}\left(\frac{M_w}{M_d}\right)
=-\frac{2h_m}{R(t)}[U_s-C_a(t)].
\end{equation}
式\eqref{eq_mass_shrink}由材料域方程积分并使用表面条件得到，可不指定未知$s_0$而检验归一化水分收支。关闭边界交换时，均匀初始含水率在纯收缩下保持不变；令$R$为常数时，模型退化为采用同一附录4物性的固定域方程。两种退化测试均与现有程序结果一致。
\FloatBarrier
